\ELhold

\ELsection{سرفصل‌ها}{سرفصل‌ها}

\begin{ELunit}

\ELfigure{m01-course-outline}{}

\end{ELunit}

\ELhold

\ELsection{آنالیز برداری}{آنالیز برداری}

\begin{ELunit}

\begin{itemize}
\tightlist
\item
  \textbf{تحلیل یا آنالیز برداری}، یک ابزار ریاضی است که به کمک آن بیان و درک مفاهیم الکترومغناطیس ساده می‌شود.
\item
  کمیت‌ها در الکترومغناطیس به دو دسته زیر تقسیم می‌شوند.

  \begin{itemize}
  \tightlist
  \item
    \textbf{عددی (\lr{scalar}):} کمیتی است که به طور کامل با اندازه‌اش (مثبت یا منفی) مشخص می‌شود. مانند: بار، جریان الکتریکی، انرژی و\ldots{}
  \item
    \textbf{برداری (\lr{vector}):} کمیتی که با اندازه و جهت توصیف می‌شود. مانند شدت میدان الکتریکی و مغناطیسی و\ldots{}
  \end{itemize}
\item
  هر دوی کمیت‌های فوق می‌توانند تابعی از مکان و زمان باشند.
\end{itemize}

\end{ELunit}

\ELnextnum{1-1}

\begin{ELexample}{}

\begin{itemize}
\tightlist
\item
  به نظر شما کمیت‌های زیر از چه نوعی هستند؟

  \begin{itemize}
  \tightlist
  \item
    دما
  \item
    توان
  \item
    نیرو
  \item
    اختلاف پتانسیل
  \item
    سرعت
  \end{itemize}
\item
  به غیر از موارد فوق، چه کمیت‌های اسکالر و برداری دیگری را در حوزه فیزیک می‌توانید نام ببرید؟
\end{itemize}

\end{ELexample}

\begin{ELunit}

\ELfigure{m01-vector-analysis-map}{}

\begin{itemize}
\tightlist
\item
  چگونه می‌توان دو کمیت برداری را به صورت ترسیمی با هم جمع یا از هم کم کرد؟
\item
  چگونه می‌توان بردارها را در هم ضرب کرد؟ مفهوم این ضرب چیست؟ کاربردهای آن کجاست؟
\end{itemize}

\ELdisp{\vect{A}\ELbr =2\uvec{x}+3\uvec{y}+\uvec{z}\qquad\ELbr  \vect{B}\ELbr =\uvec{x}-2\uvec{y}+3\uvec{z}\qquad\ELbr  \vect{C}\ELbr =-\uvec{x}+4\uvec{y}-\uvec{z}}

\begin{itemize}
\tightlist
\item
  اندازه تصویر بردار \ELim{\vect{A}} بر بردار \ELim{\vect{B}} را بدست آورید.
\item
  زاویه بین بردارهای \ELim{\vect{A}} و \ELim{\vect{B}} چقدر است؟
\item
  مساحت مثلث تشکیل شده توسط دو بردار \ELim{\vect{A}} و \ELim{\vect{B}} چقدر می‌شود؟
\item
  بردار واحد عمود بر صفحه شامل دو بردار \ELim{\vect{A}} و \ELim{\vect{B}} را بدست آورید.
\item
  حجم متوازی‌السطوح حاصل از ۳ بردار \ELim{\vect{A}}، \ELim{\vect{B}} و \ELim{\vect{C}} چقدر است؟
\end{itemize}

\end{ELunit}

\ELhold

\ELsection{جبر برداری}{جبر برداری}

\begin{ELunit}

\begin{itemize}
\tightlist
\item
  هر برداری دارای اندازه و جهت است.
\end{itemize}

\ELdisp{\vect{A}\ELbr =\uvec{A}A,\qquad\ELbr  \uvec{A}\ELbr =\frac{\vect{A}}{\abs{\vect{A}}}\ELbr =\frac{\vect{A}}{A}\quad\ELbr (\ELt{بردار یکه})}

\ELfigure{m01-vector-magnitude}{}

\end{ELunit}

\begin{ELremark}{}

برای هر کمیت برداری باید هم اندازه و هم جهت معلوم باشد. عدم مشخص کردن جهت یک کمیت برداری، به معنای نادیده گرفتن بخشی از اطلاعات آن کمیت است.

\end{ELremark}

\ELhold

\ELsection{جمع بردارها}{جمع بردارها}

\begin{ELunit}

\begin{itemize}
\tightlist
\item
  جمع دو بردار به صورت ترسیمی

  \begin{itemize}
  \tightlist
  \item
    قاعده \textbf{متوازی‌الاضلاع}
  \item
    قاعده \textbf{ابتدا به انتها}
  \end{itemize}
\end{itemize}

\ELfigure{m01-vector-sum}{}

\begin{itemize}
\tightlist
\item
  قوانین حاکم بر جمع بردارها

  \begin{itemize}
  \item
    \textbf{قانون جابجایی}
    \ELdisp{\vect{A}+\vect{B}\ELbr =\vect{B}+\vect{A}}
  \item
    \textbf{قانون انجمنی}
    \ELdisp{\vect{A}+(\vect{B}+\vect{C})\ELbr =(\vect{A}+\vect{B})+\vect{C}}
  \end{itemize}
\end{itemize}

\end{ELunit}

\ELhold

\ELsection{تفریق بردارها}{تفریق بردارها}

\begin{ELunit}

\begin{itemize}
\tightlist
\item
  تفریق برداری
\end{itemize}

\ELdisp{\vect{A}-\vect{B}\ELbr =\vect{A}+(-\vect{B})}

\ELfigure{m01-vector-difference}{}

\end{ELunit}

\ELhold

\ELsection{ضرب بردارها}{ضرب بردارها}

\begin{ELunit}

\ELfigure{m01-products-map}{}

\begin{itemize}
\tightlist
\item
  ضرب یک اسکالر در بردار:
\end{itemize}

\ELdisp{k\vect{A}\ELbr =\uvec{A}(kA)}

\end{ELunit}

\ELhold

\ELsection{ضرب داخلی یا عددی}{ضرب داخلی یا عددی}

\begin{ELdefinition}{تعریف ریاضی}

\ELdisp{\vect{A}\cdot\vect{B}\triangleq AB\cos\theta_{AB}}

\begin{itemize}
\tightlist
\item
  برابر است با حاصلضرب اندازه یک بردار در تصویر بردار دیگر بر بردار اول.
\end{itemize}

\end{ELdefinition}

\begin{ELunit}

\ELfigure{m01-dot-product}{\ELim{\theta_{AB}}: زاویه کوچکتر بین دو بردار هنگامیکه ابتدای دو بردار
به هم متصل می‌شود}

\begin{itemize}
\tightlist
\item
  ویژگی‌های ضرب داخلی

  \begin{itemize}
  \item
    نتیجه ضرب داخلی دو بردار، یک کمیت اسکالر است.
  \item
    کمتر یا مساوی حاصلضرب اندازه‌های آن‌هاست.
  \item
    بسته به زاویه بین آن‌ها می‌تواند کمیتی مثبت یا منفی باشد.
  \item
    اگر بردارها عمود بر هم باشند، مساوی صفر است.
  \item
    جابجاپذیر و توزیع‌پذیر است.
    \ELdisp{\vect{A}\cdot\vect{B}\ELbr =\vect{B}\cdot\vect{A},\qquad\ELbr  \vect{A}\cdot(\vect{B}+\vect{C})\ELbr =\vect{A}\cdot\vect{B}+\vect{A}\cdot\vect{C}}
  \item
    ضرب داخلی یک بردار در خودش
    \ELdisp{\vect{A}\cdot\vect{A}\ELbr =A^2,\qquad\ELbr  A\ELbr =+\sqrt{\vect{A}\cdot\vect{A}}}
  \end{itemize}
\end{itemize}

\end{ELunit}

\ELnextnum{1-2}

\begin{ELexample}{}

با استفاده از مفهوم ضرب داخلی تصویر بردار \ELim{\vect{A}} بر \ELim{\vect{B}} را بدست آورید.

\ELfigure{m01-ex2-statement}{}

\begin{ELsolution}

\ELfigure{m01-ex2-projection}{}

اندازه تصویر بردار \ELim{\vect{A}} بر \ELim{\vect{B}}:
\ELdisp{A\cos\theta_{AB}\ELbr =\frac{\vect{A}\cdot\vect{B}}{B}}
تصویر بردار \ELim{\vect{A}} بر \ELim{\vect{B}}:
\ELdisp{\frac{\vect{A}\cdot\vect{B}}{B}\uvec{B}\ELbr =\left(\frac{\vect{A}\cdot\vect{B}}{B}\right)\frac{\vect{B}}{B}\ELbr =\frac{\vect{A}\cdot\vect{B}}{B^2}\vect{B}}

\end{ELsolution}

\end{ELexample}

\ELnextnum{1-3}

\begin{ELexample}{}

قانون کسینوس‌ها در یک مثلث را با استفاده از مفهوم ضرب داخلی اثبات کنید.
\ELdisp{C\ELbr =\sqrt{A^2+B^2-2AB\cos\alpha}}

\ELfigure{m01-ex3-triangle}{}

\begin{ELsolution}

\ELfigure{m01-ex3-solution}{}

\ELdisp{\vect{C}\ELbr =\vect{A}+\vect{B}}
\ELdisp{\begin{aligned}C^2=\vect{C}\cdot\vect{C}&=(\vect{A}+\vect{B})\cdot(\vect{A}+\vect{B})\\&=\vect{A}\cdot\vect{A}+\vect{B}\cdot\vect{B}+2\vect{A}\cdot\vect{B}\\&=A^2+B^2+2AB\cos\theta_{AB}\end{aligned}}
\ELdisp{\cos\theta_{AB}\ELbr =\cos(180^\circ-\alpha)\ELbr =-\cos\alpha}
\ELdisp{C^2\ELbr =A^2+B^2-2AB\cos\alpha}

\end{ELsolution}

\end{ELexample}

\ELhold

\ELsection{ضرب خارجی یا برداری}{ضرب خارجی یا برداری}

\begin{ELdefinition}{}

\ELdisp{\vect{A}\times\vect{B}\triangleq\uvec{n}\abs{AB\sin\theta_{AB}}}

\end{ELdefinition}

\begin{ELunit}

\begin{itemize}
\tightlist
\item
  اندازه این حاصلضرب برابر با مساحت متوازی‌الاضلاع ساخته شده توسط بردارهای \ELim{\vect{A}} و \ELim{\vect{B}} و جهت آن عمود بر هر دو بردار است.
\end{itemize}

\ELfigure{m01-cross-product}{}

\end{ELunit}

\begin{ELjoin}

\begin{itemize}
\tightlist
\item
  ویژگی‌های ضرب خارجی

  \begin{itemize}
  \item
    نتیجه ضرب خارجی دو بردار یک کمیت برداری است.
  \item
    جابجاپذیر نیست
    \ELdisp{\vect{B}\times\vect{A}\ELbr =-\vect{A}\times\vect{B}}
  \item
    توزیع‌پذیر است
    \ELdisp{\vect{A}\times(\vect{B}+\vect{C})\ELbr =\vect{A}\times\vect{B}+\vect{A}\times\vect{C}}
  \item
    انجمن‌پذیر نیست
    \ELdisp{\vect{A}\times(\vect{B}\times\vect{C})\neq(\vect{A}\times\vect{B})\times\vect{C}}
  \end{itemize}
\end{itemize}

\ELnextnum{1-4}

\begin{ELexample}{}

قانون سینوس‌ها در یک مثلث را با استفاده از مفهوم ضرب خارجی اثبات کنید.
\ELdisp{\frac{\sin A}{a}\ELbr =\frac{\sin B}{b}\ELbr =\frac{\sin C}{c}}

\ELfigure{m01-ex4-triangle}{}

\begin{ELsolution}

مساحت مثلث:
\ELdisp{\abs*{\tfrac12\vect{a}\times\vect{b}}\ELbr =\abs*{\tfrac12\vect{b}\times\vect{c}}\ELbr =\abs*{\tfrac12\vect{c}\times\vect{a}}}
\ELdisp{ab\sin C\ELbr =bc\sin A\ELbr =ca\sin B}
\ELdisp{\frac{\sin A}{a}\ELbr =\frac{\sin B}{b}\ELbr =\frac{\sin C}{c}}

\end{ELsolution}

\end{ELexample}

\end{ELjoin}

\ELhold

\ELsection{ضرب سه بردار}{ضرب سه بردار}

\begin{ELunit}

\begin{itemize}
\tightlist
\item
  \textbf{ضرب سه‌گانه عددی}
\end{itemize}

\ELdisp{\vect{A}\cdot(\vect{B}\times\vect{C})\ELbr =\vect{B}\cdot(\vect{C}\times\vect{A})\ELbr =\vect{C}\cdot(\vect{A}\times\vect{B})}

\begin{itemize}
\tightlist
\item
  دارای اندازه‌ای برابر با حجم متوازی‌السطوح تشکیل یافته از سه بردار \ELim{\vect{A}}، \ELim{\vect{B}} و \ELim{\vect{C}}.

  \begin{itemize}
  \tightlist
  \item
    مساحت قاعده: \ELim{\abs{\vect{B}\times\vect{C}}=\abs{BC\sin\theta_1}}
  \item
    اندازه ارتفاع: \ELim{\abs{A\cos\theta_2}}
  \item
    حجم متوازی‌السطوح: \ELim{\abs{ABC\sin\theta_1\cos\theta_2}}
  \end{itemize}
\end{itemize}

\ELfigure{m01-triple-product}{}

\begin{itemize}
\tightlist
\item
  \textbf{ضرب سه‌گانه برداری}

  \begin{itemize}
  \tightlist
  \item
    معروف به قاعده \lr{back-cab}
  \end{itemize}
\end{itemize}

\ELdisp{\vect{A}\times(\vect{B}\times\vect{C})\ELbr =\vect{B}(\vect{A}\cdot\vect{C})-\vect{C}(\vect{A}\cdot\vect{B})}

\end{ELunit}

\ELnextnum{1-5}

\begin{ELexample}{}

عبارت \ELim{(\vect{A}\times\vect{B})\times\vect{C}} با کدامیک از عبارت‌های زیر برابر است؟

\begin{itemize}
\tightlist
\item
  \ELim{\vect{B}(\vect{A}\cdot\vect{C})-\vect{C}(\vect{A}\cdot\vect{B})}
\item
  \ELim{-\vect{B}(\vect{A}\cdot\vect{C})+\vect{C}(\vect{A}\cdot\vect{B})}
\item
  \ELim{-\vect{A}(\vect{C}\cdot\vect{B})+\vect{B}(\vect{C}\cdot\vect{A})}
\item
  \ELim{\vect{A}(\vect{C}\cdot\vect{B})-\vect{B}(\vect{C}\cdot\vect{A})}
\end{itemize}

\begin{ELsolution}

\ELdisp{\begin{aligned}(\vect{A}\times\vect{B})\times\vect{C}&=-\vect{C}\times(\vect{A}\times\vect{B})\\&=-\vect{A}(\vect{C}\cdot\vect{B})+\vect{B}(\vect{C}\cdot\vect{A})\end{aligned}}

\end{ELsolution}

\end{ELexample}

\ELhold

\ELsection{دستگاه‌های مختصات متعامد}{دستگاه‌های مختصات متعامد}

\begin{ELunit}

\ELfigure{m01-vector-analysis-map-coord}{}

\begin{itemize}
\item
  صفحات اصلی در دستگاه‌های مختصات مختلف کدامند؟
\item
  بردارهای یکه در دستگاه‌های مختصات مختلف چگونه تعریف می‌شوند؟
\item
  نحوه محاسبه بردار مکان یک نقطه و بردار فاصله بین دو نقطه در دستگاه‌های مختصات مختلف چگونه است؟
\item
  طول، سطح و حجم دیفرانسیلی را در دستگاه‌های مختصات مختلف، جهت استفاده در انتگرال‌گیری، تعیین کنید.
\item
  نمایش بردار \ELim{\vect{A}=\uvec{r}(3\cos\phi)-\uvec{\phi}2r+\uvec{z}5} در مختصات کارتزین را بدست آورید.
\item
  موقعیت نقطه \ELim{P} در مختصات کروی \ELim{(8,120^\circ,330^\circ)} است. موقعیت این نقطه را در مختصات استوانه‌ای و کارتزین بیان کنید.
\item
  برای توصیف تغییرات مکانی کمیت‌های فیزیکی باید بتوانیم تمام نقاط فضا را به صورت یکتا و مناسب توصیف کنیم.

  \begin{itemize}
  \tightlist
  \item
    این امر مستلزم بکارگیری یک دستگاه مختصات مناسب است.
  \end{itemize}
\item
  قوانین الکترومغناطیسی به دستگاه مختصات وابسته نیستند.
\item
  جهت سادگی محاسبات از دستگاه مختصاتی که مناسب هندسه ساختار مسئله باشد، استفاده می‌شود.
\item
  در فضای سه بعدی یک نقطه می‌تواند در محل تقاطع سه سطح قرار گیرد.

  \begin{itemize}
  \tightlist
  \item
    اگر این سه سطح دو به دو بر هم عمود باشند، یک دستگاه مختصات متعامد خواهیم داشت.
  \end{itemize}
\item
  سه دستگاه متعامد مورد استفاده در این درس

  \begin{itemize}
  \tightlist
  \item
    مختصات کارتزین یا دکارتی یا قائم
  \item
    مختصات استوانه‌ای
  \item
    مختصات کروی
  \end{itemize}
\end{itemize}

\end{ELunit}

\ELhold

\ELsection{دستگاه مختصات کارتزین}{دستگاه مختصات کارتزین}

\begin{ELunit}

\ELfigure{m01-cartesian-planes}{}

\begin{itemize}
\item
  هر نقطه \ELim{P(x_1,y_1,z_1)} محل تقاطع سه صفحه \ELim{x=x_1}، \ELim{y=y_1} و \ELim{z=z_1} است.
\item
  بردار یکه \ELim{\uvec{x}} \ELim{\ELto} بردار واحد عمود بر صفحه \ELim{x}-ثابت و در جهت مثبت
\item
  بردار یکه \ELim{\uvec{y}} \ELim{\ELto} بردار واحد عمود بر صفحه \ELim{y}-ثابت و در جهت مثبت
\item
  بردار یکه \ELim{\uvec{z}} \ELim{\ELto} بردار واحد عمود بر صفحه \ELim{z}-ثابت و در جهت مثبت
\item
  ضرب خارجی و داخلی بردارهای یکه
\end{itemize}

\ELdisp{\uvec{x}\times\uvec{y}\ELbr =\uvec{z},\qquad\ELbr  \uvec{y}\times\uvec{z}\ELbr =\uvec{x},\qquad\ELbr  \uvec{z}\times\uvec{x}\ELbr =\uvec{y}}

\ELfigure{m01-cartesian-cycle}{}

\ELdisp{\uvec{x}\cdot\uvec{y}\ELbr =\uvec{y}\cdot\uvec{z}\ELbr =\uvec{z}\cdot\uvec{x}\ELbr =0}
\ELdisp{\uvec{x}\cdot\uvec{x}\ELbr =\uvec{y}\cdot\uvec{y}\ELbr =\uvec{z}\cdot\uvec{z}\ELbr =1}

\end{ELunit}

\begin{ELjoin}

\begin{itemize}
\tightlist
\item
  نمایش بردار دلخواه \ELim{\vect{A}} در مختصات کارتزین
\end{itemize}

\begin{ELimportant}{}

\ELdisp{\vect{A}\ELbr =\uvec{x}A_x+\uvec{y}A_y+\uvec{z}A_z}
\ELdisp{A\ELbr =\sqrt{A_x^2+A_y^2+A_z^2}}

\end{ELimportant}

\end{ELjoin}

\begin{ELunit}

\begin{itemize}
\tightlist
\item
  \textbf{بردار مکان}

  \begin{itemize}
  \tightlist
  \item
    بردار مکان نقطه \ELim{P}، فاصله جهت‌دار از مبدأ به \ELim{P} است.
  \end{itemize}
\end{itemize}

\ELdisp{\vect{R}_1\ELbr =x_1\uvec{x}+y_1\uvec{y}+z_1\uvec{z},\qquad\ELbr  \vect{R}_2\ELbr =x_2\uvec{x}+y_2\uvec{y}+z_2\uvec{z}}

\begin{itemize}
\tightlist
\item
  \textbf{بردار فاصله}

  \begin{itemize}
  \tightlist
  \item
    بردار جابجایی از یک نقطه به نقطه دیگر است.
  \end{itemize}
\end{itemize}

\ELdisp{\vect{R}_{12}\ELbr =\vect{R}_2-\vect{R}_1\ELbr =(x_2-x_1)\uvec{x}+(y_2-y_1)\uvec{y}+(z_2-z_1)\uvec{z}}

\ELfigure{m01-position-vectors}{}

\ELdisp{\vect{A}\ELbr =A_x\uvec{x}+A_y\uvec{y}+A_z\uvec{z},\qquad\ELbr  \vect{B}\ELbr =B_x\uvec{x}+B_y\uvec{y}+B_z\uvec{z}}

\end{ELunit}

\begin{ELjoin}

\begin{itemize}
\tightlist
\item
  ضرب داخلی دو بردار \ELim{\vect{A}} و \ELim{\vect{B}}
\end{itemize}

\begin{ELimportant}{}

\ELdisp{\vect{A}\cdot\vect{B}\ELbr =A_xB_x+A_yB_y+A_zB_z}

\end{ELimportant}

\end{ELjoin}

\begin{ELjoin}

\begin{itemize}
\tightlist
\item
  ضرب خارجی دو بردار \ELim{\vect{A}} و \ELim{\vect{B}}
\end{itemize}

\ELdisp{\begin{aligned}\vect{A}\times\vect{B}={}&A_xB_x\uvec{x}\times\uvec{x}+A_xB_y\uvec{x}\times\uvec{y}+A_xB_z\uvec{x}\times\uvec{z}\\&+A_yB_x\uvec{y}\times\uvec{x}+A_yB_y\uvec{y}\times\uvec{y}+A_yB_z\uvec{y}\times\uvec{z}\\&+A_zB_x\uvec{z}\times\uvec{x}+A_zB_y\uvec{z}\times\uvec{y}+A_zB_z\uvec{z}\times\uvec{z}\\={}&(A_yB_z-A_zB_y)\uvec{x}+(A_zB_x-A_xB_z)\uvec{y}+(A_xB_y-A_yB_x)\uvec{z}\end{aligned}}

\begin{ELimportant}{}

\ELdisp{\vect{A}\times\vect{B}\ELbr =\begin{vmatrix}\uvec{x}&\uvec{y}&\uvec{z}\\A_x&A_y&A_z\\B_x&B_y&B_z\end{vmatrix}}

\end{ELimportant}

\end{ELjoin}

\begin{ELunit}

\begin{itemize}
\tightlist
\item
  طول دیفرانسیلی \ELim{\ELto} \textbf{بردار}
\end{itemize}

\ELdisp{\dif\vect{\ell}\ELbr =\uvec{x}\,dx+\uvec{y}\,dy+\uvec{z}\,dz}

\begin{itemize}
\tightlist
\item
  سطح دیفرانسیلی \ELim{\ELto} \textbf{بردار}
\end{itemize}

\ELdisp{d\vect{S}\ELbr =dS\,\uvec{n}}
\ELdisp{d\vect{S}\ELbr =dy\,dz\,\uvec{x},\qquad\ELbr  d\vect{S}\ELbr =dx\,dz\,\uvec{y},\qquad\ELbr  d\vect{S}\ELbr =dx\,dy\,\uvec{z}}

\begin{itemize}
\tightlist
\item
  حجم دیفرانسیلی \ELim{\ELto} \textbf{اسکالر}
\end{itemize}

\ELdisp{dv\ELbr =dx\,dy\,dz}

\ELfigure{m01-cartesian-differentials}{}

\end{ELunit}

\ELnextnum{1-6}

\begin{ELexample}{}

الف) برداری که نقطه ابتدای آن \ELim{P_1(1,3,2)} و نقطه انتهای آن \ELim{P_2(3,-2,4)} است را بدست آورید.

ب) طول این بردار چقدر است؟

\begin{ELsolution}

\ELfigure{m01-ex6}{}

الف)
\ELdisp{\begin{aligned}\overrightarrow{P_1P_2}&=\overrightarrow{OP_2}-\overrightarrow{OP_1}\\&=(\uvec{x}3-\uvec{y}2+\uvec{z}4)-(\uvec{x}+\uvec{y}3+\uvec{z}2)\\&=\uvec{x}2-\uvec{y}5+\uvec{z}2\end{aligned}}
ب)
\ELdisp{P_1P_2\ELbr =\abs{\overrightarrow{P_1P_2}}\ELbr =\sqrt{2^2+(-5)^2+2^2}\ELbr =\sqrt{33}}

\end{ELsolution}

\end{ELexample}

\ELnextnum{1-7}

\begin{ELexample}{}

با فرض اینکه \ELim{\vect{A}=\uvec{x}5-\uvec{y}2+\uvec{z}}، آنگاه بردار یکه \ELim{\vect{B}} را به گونه‌ای بدست آورید که

\begin{itemize}
\tightlist
\item
  الف. \ELim{\vect{B}\parallel\vect{A}}
\item
  ب. \ELim{\vect{B}} در صفحه \ELim{xy} بوده و عمود بر \ELim{\vect{A}} باشد
\end{itemize}

\begin{ELsolution}

الف.
\ELdisp{\vect{B}\ELbr =\frac{\vect{A}}{A}\ELbr =\frac{\uvec{x}5-\uvec{y}2+\uvec{z}}{\sqrt{5^2+2^2+1^2}}\ELbr =\frac{\uvec{x}5-\uvec{y}2+\uvec{z}}{\sqrt{30}}}
ب.
\ELdisp{\begin{cases}B=\sqrt{B_x^2+B_y^2+B_z^2}=1\\B_z=0\\\vect{B}\cdot\vect{A}=5B_x-2B_y=0\end{cases}\quad\ELbr \ELbr \Rightarrow\quad\ELbr \begin{cases}B_x=\dfrac{2}{\sqrt{29}}\\\noalign{\vskip 2mm}B_y=\dfrac{5}{\sqrt{29}}\end{cases}}

\end{ELsolution}

\end{ELexample}

\ELhold

\ELsection{دستگاه مختصات استوانه‌ای}{دستگاه مختصات استوانه‌ای}

\begin{ELunit}

\ELfigure{m01-cylindrical-surfaces}{}

\begin{itemize}
\tightlist
\item
  هر نقطه \ELim{P(r_1,\phi_1,z_1)} محل تقاطع استوانه \ELim{r=r_1}، نیم‌صفحه \ELim{\phi=\phi_1} و صفحه \ELim{z=z_1} است.
\item
  بردار یکه \ELim{\uvec{r}} \ELim{\ELto} بردار واحد عمود بر استوانه \ELim{r}-ثابت و در جهت مثبت
\item
  بردار یکه \ELim{\uvec{\phi}} \ELim{\ELto} بردار واحد عمود بر نیم‌صفحه \ELim{\phi}-ثابت و در جهت مثبت
\item
  بردار یکه \ELim{\uvec{z}} \ELim{\ELto} بردار واحد عمود بر صفحه \ELim{z}-ثابت و در جهت مثبت
\end{itemize}

\end{ELunit}

\begin{ELremark}{}

جهت بردارهای \ELim{\uvec{r}} و \ELim{\uvec{\phi}} در نقاط مختلف فضا تغییر می‌کند. این نکته در هنگام انتگرال‌گیری باید مورد دقت قرار گیرد.

\end{ELremark}

\begin{ELunit}

\begin{itemize}
\tightlist
\item
  ضرب خارجی و داخلی بردارهای یکه
\end{itemize}

\ELdisp{\uvec{r}\times\uvec{\phi}\ELbr =\uvec{z},\qquad\ELbr  \uvec{\phi}\times\uvec{z}\ELbr =\uvec{r},\qquad\ELbr  \uvec{z}\times\uvec{r}\ELbr =\uvec{\phi}}

\ELfigure{m01-cylindrical-cycle}{}

\ELdisp{\uvec{r}\cdot\uvec{r}\ELbr =\uvec{\phi}\cdot\uvec{\phi}\ELbr =\uvec{z}\cdot\uvec{z}\ELbr =1}
\ELdisp{\uvec{r}\cdot\uvec{\phi}\ELbr =\uvec{\phi}\cdot\uvec{z}\ELbr =\uvec{z}\cdot\uvec{r}\ELbr =0}

\end{ELunit}

\begin{ELjoin}

\begin{itemize}
\tightlist
\item
  نمایش بردار دلخواه \ELim{\vect{A}} در مختصات استوانه‌ای
\end{itemize}

\begin{ELimportant}{}

\ELdisp{\vect{A}\ELbr =\uvec{r}A_r+\uvec{\phi}A_\phi+\uvec{z}A_z}
\ELdisp{A\ELbr =\sqrt{A_r^2+A_\phi^2+A_z^2}}

\end{ELimportant}

\end{ELjoin}

\begin{ELunit}

\begin{itemize}
\tightlist
\item
  بردار مکان
\end{itemize}

\ELfigure{m01-cylindrical-position}{}

\ELdisp{\vect{R}\ELbr =r\uvec{r}+z\uvec{z}}

\end{ELunit}

\ELnextnum{1-8}

\begin{ELexample}{}

مختصات استوانه‌ای نقطه دلخواه \ELim{P} به صورت \ELim{(r,\phi,0)} است. بردار واحد از نقطه \ELim{z=h} روی محور \ELim{z} به سمت نقطه \ELim{P} را بیابید.

\begin{ELsolution}

\ELfigure{m01-ex8}{}

\ELdisp{\overrightarrow{QP}\ELbr =\overrightarrow{OP}-\overrightarrow{OQ}\ELbr =(\uvec{r}r)-(\uvec{z}h)}
\ELdisp{\uvec{QP}\ELbr =\frac{\overrightarrow{QP}}{\abs{\overrightarrow{QP}}}\ELbr =\frac{1}{\sqrt{r^2+h^2}}(\uvec{r}r-\uvec{z}h)}

\end{ELsolution}

\end{ELexample}

\begin{ELunit}

\begin{itemize}
\tightlist
\item
  طول دیفرانسیلی \ELim{\ELto} \textbf{بردار}
\end{itemize}

\ELdisp{\dif\vect{\ell}\ELbr =\uvec{r}\,dr+\uvec{\phi}\,r\,d\phi+\uvec{z}\,dz}

\begin{itemize}
\tightlist
\item
  سطح دیفرانسیلی \ELim{\ELto} \textbf{بردار}
\end{itemize}

\ELdisp{d\vect{S}\ELbr =r\,d\phi\,dz\,\uvec{r},\qquad\ELbr  d\vect{S}\ELbr =dr\,dz\,\uvec{\phi},\qquad\ELbr  d\vect{S}\ELbr =r\,dr\,d\phi\,\uvec{z}}

\begin{itemize}
\tightlist
\item
  حجم دیفرانسیلی \ELim{\ELto} \textbf{اسکالر}
\end{itemize}

\ELdisp{dv\ELbr =r\,dr\,d\phi\,dz}

\ELfigure{m01-cylindrical-differentials}{}

\begin{itemize}
\tightlist
\item
  تبدیل نمایش بردار از مختصات استوانه‌ای به کارتزین
\end{itemize}

\ELdisp{\vect{A}\ELbr =\uvec{r}A_r+\uvec{\phi}A_\phi+\uvec{z}A_z\quad\ELbr \ELbr \Longrightarrow\quad\ELbr \vect{A}\ELbr =\uvec{x}A_x+\uvec{y}A_y+\uvec{z}A_z}

\ELfigure{m01-cyl-cart-units}{}

\ELdisp{\uvec{r}\ELbr =\cos\phi\,\uvec{x}+\sin\phi\,\uvec{y},\qquad\ELbr  \uvec{\phi}\ELbr =-\sin\phi\,\uvec{x}+\cos\phi\,\uvec{y},\qquad\ELbr  \uvec{z}\ELbr =\uvec{z}}

\ELdisp{\begin{aligned}\vect{A}&=(\cos\phi\,\uvec{x}+\sin\phi\,\uvec{y})A_r\\&\quad+(-\sin\phi\,\uvec{x}+\cos\phi\,\uvec{y})A_\phi\\&\quad+\uvec{z}A_z\end{aligned}\quad\ELbr \ELbr \Longrightarrow\quad\ELbr \begin{aligned}\vect{A}&=\uvec{x}(\cos\phi\,A_r-\sin\phi\,A_\phi)\\&\quad+\uvec{y}(\sin\phi\,A_r+\cos\phi\,A_\phi)\\&\quad+\uvec{z}A_z\end{aligned}}

\begin{itemize}
\tightlist
\item
  تبدیل نمایش بردار از مختصات کارتزین به استوانه‌ای
\end{itemize}

\ELdisp{\vect{A}\ELbr =\uvec{x}A_x+\uvec{y}A_y+\uvec{z}A_z\quad\ELbr \ELbr \Longrightarrow\quad\ELbr \vect{A}\ELbr =\uvec{r}A_r+\uvec{\phi}A_\phi+\uvec{z}A_z}

\ELdisp{\uvec{x}\ELbr =\cos\phi\,\uvec{r}-\sin\phi\,\uvec{\phi},\qquad\ELbr  \uvec{y}\ELbr =\sin\phi\,\uvec{r}+\cos\phi\,\uvec{\phi},\qquad\ELbr  \uvec{z}\ELbr =\uvec{z}}

\ELdisp{\begin{aligned}\vect{A}&=(\cos\phi\,\uvec{r}-\sin\phi\,\uvec{\phi})A_x\\&\quad+(\sin\phi\,\uvec{r}+\cos\phi\,\uvec{\phi})A_y\\&\quad+\uvec{z}A_z\end{aligned}\quad\ELbr \ELbr \Longrightarrow\quad\ELbr \begin{aligned}\vect{A}&=\uvec{r}(\cos\phi\,A_x+\sin\phi\,A_y)\\&\quad+\uvec{\phi}(-\sin\phi\,A_x+\cos\phi\,A_y)\\&\quad+\uvec{z}A_z\end{aligned}}

\end{ELunit}

\begin{ELjoin}

\begin{itemize}
\tightlist
\item
  تبدیل نمایش بردار از مختصات استوانه‌ای به کارتزین
\end{itemize}

\begin{ELimportant}{}

\ELdisp{\begin{bmatrix}A_x\\A_y\\A_z\end{bmatrix}\ELbr =\begin{bmatrix}\cos\phi&-\sin\phi&0\\\sin\phi&\cos\phi&0\\0&0&1\end{bmatrix}\begin{bmatrix}A_r\\A_\phi\\A_z\end{bmatrix}}

\end{ELimportant}

\end{ELjoin}

\begin{ELjoin}

\begin{itemize}
\tightlist
\item
  تبدیل نمایش بردار از مختصات کارتزین به استوانه‌ای
\end{itemize}

\begin{ELimportant}{}

\ELdisp{\begin{bmatrix}A_r\\A_\phi\\A_z\end{bmatrix}\ELbr =\begin{bmatrix}\cos\phi&\sin\phi&0\\-\sin\phi&\cos\phi&0\\0&0&1\end{bmatrix}\begin{bmatrix}A_x\\A_y\\A_z\end{bmatrix}}

\end{ELimportant}

\end{ELjoin}

\begin{ELunit}

\begin{itemize}
\tightlist
\item
  تبدیل از مختصات استوانه‌ای به کارتزین
\end{itemize}

\ELdisp{x\ELbr =r\cos\phi,\qquad\ELbr  y\ELbr =r\sin\phi,\qquad\ELbr  z\ELbr =z}

\begin{itemize}
\tightlist
\item
  تبدیل از مختصات کارتزین به استوانه‌ای
\end{itemize}

\ELdisp{r\ELbr =\sqrt{x^2+y^2},\qquad\ELbr  \phi\ELbr =\tan^{-1}\frac{y}{x},\qquad\ELbr  z\ELbr =z}

\ELfigure{m01-cyl-cart-coords}{}

\end{ELunit}

\ELhold

\ELsection{دستگاه مختصات کروی}{دستگاه مختصات کروی}

\begin{ELunit}

\ELfigure{m01-spherical-surfaces}{}

\begin{itemize}
\tightlist
\item
  هر نقطه \ELim{P(R_1,\theta_1,\phi_1)} محل تقاطع کره \ELim{R=R_1}، مخروط \ELim{\theta=\theta_1} و نیم‌صفحه \ELim{\phi=\phi_1}
\item
  بردار یکه \ELim{\uvec{R}} \ELim{\ELto} بردار واحد عمود بر کره \ELim{R}-ثابت و در جهت مثبت
\item
  بردار یکه \ELim{\uvec{\theta}} \ELim{\ELto} بردار واحد عمود بر مخروط \ELim{\theta}-ثابت و در جهت مثبت
\item
  بردار یکه \ELim{\uvec{\phi}} \ELim{\ELto} بردار واحد عمود بر نیم‌صفحه \ELim{\phi}-ثابت و در جهت مثبت
\end{itemize}

\end{ELunit}

\begin{ELremark}{}

جهت بردارهای \ELim{\uvec{R}}، \ELim{\uvec{\theta}} و \ELim{\uvec{\phi}} در نقاط مختلف فضا تغییر می‌کند. این نکته در هنگام انتگرال‌گیری باید مورد دقت قرار گیرد.

\end{ELremark}

\begin{ELunit}

\begin{itemize}
\tightlist
\item
  ضرب خارجی و داخلی بردارهای یکه
\end{itemize}

\ELdisp{\uvec{R}\times\uvec{\theta}\ELbr =\uvec{\phi},\qquad\ELbr  \uvec{\theta}\times\uvec{\phi}\ELbr =\uvec{R},\qquad\ELbr  \uvec{\phi}\times\uvec{R}\ELbr =\uvec{\theta}}

\ELfigure{m01-spherical-cycle}{}

\ELdisp{\uvec{R}\cdot\uvec{R}\ELbr =\uvec{\theta}\cdot\uvec{\theta}\ELbr =\uvec{\phi}\cdot\uvec{\phi}\ELbr =1}
\ELdisp{\uvec{R}\cdot\uvec{\theta}\ELbr =\uvec{\theta}\cdot\uvec{\phi}\ELbr =\uvec{\phi}\cdot\uvec{R}\ELbr =0}

\end{ELunit}

\begin{ELjoin}

\begin{itemize}
\tightlist
\item
  نمایش بردار دلخواه \ELim{\vect{A}} در مختصات کروی
\end{itemize}

\begin{ELimportant}{}

\ELdisp{\vect{A}\ELbr =\uvec{R}A_R+\uvec{\theta}A_\theta+\uvec{\phi}A_\phi}
\ELdisp{A\ELbr =\sqrt{A_R^2+A_\theta^2+A_\phi^2}}

\end{ELimportant}

\end{ELjoin}

\begin{ELunit}

\begin{itemize}
\tightlist
\item
  بردار مکان
\end{itemize}

\ELfigure{m01-spherical-position}{}

\ELdisp{\vect{R}\ELbr =R_1\uvec{R}}

\begin{itemize}
\tightlist
\item
  طول دیفرانسیلی \ELim{\ELto} \textbf{بردار}
\end{itemize}

\ELdisp{\dif\vect{\ell}\ELbr =\uvec{R}\,dR+\uvec{\theta}\,R\,d\theta+\uvec{\phi}\,R\sin\theta\,d\phi}

\begin{itemize}
\tightlist
\item
  سطح دیفرانسیلی \ELim{\ELto} \textbf{بردار}
\end{itemize}

\ELdisp{d\vect{S}\ELbr =R^2\sin\theta\,d\theta\,d\phi\,\uvec{R},\qquad\ELbr  d\vect{S}\ELbr =R\sin\theta\,dR\,d\phi\,\uvec{\theta},\qquad\ELbr  d\vect{S}\ELbr =R\,dR\,d\theta\,\uvec{\phi}}

\begin{itemize}
\tightlist
\item
  حجم دیفرانسیلی \ELim{\ELto} \textbf{اسکالر}
\end{itemize}

\ELdisp{dv\ELbr =R^2\sin\theta\,dR\,d\theta\,d\phi}

\ELfigure{m01-spherical-differentials}{}

\begin{itemize}
\tightlist
\item
  تبدیل نمایش بردار از مختصات کروی به استوانه‌ای
\end{itemize}

\ELdisp{\vect{A}\ELbr =\uvec{R}A_R+\uvec{\theta}A_\theta+\uvec{\phi}A_\phi\quad\ELbr \ELbr \Longrightarrow\quad\ELbr \vect{A}\ELbr =\uvec{r}A_r+\uvec{\phi}A_\phi+\uvec{z}A_z}

\ELfigure{m01-sph-cyl-units}{}

\ELdisp{\uvec{R}\ELbr =\sin\theta\,\uvec{r}+\cos\theta\,\uvec{z},\qquad\ELbr  \uvec{\theta}\ELbr =\cos\theta\,\uvec{r}-\sin\theta\,\uvec{z},\qquad\ELbr  \uvec{\phi}\ELbr =\uvec{\phi}}

\ELdisp{\begin{aligned}\vect{A}&=(\sin\theta\,\uvec{r}+\cos\theta\,\uvec{z})A_R\\&\quad+(\cos\theta\,\uvec{r}-\sin\theta\,\uvec{z})A_\theta\\&\quad+\uvec{\phi}A_\phi\end{aligned}\quad\ELbr \ELbr \Longrightarrow\quad\ELbr \begin{aligned}\vect{A}&=\uvec{r}(\sin\theta\,A_R+\cos\theta\,A_\theta)\\&\quad+\uvec{\phi}A_\phi\\&\quad+\uvec{z}(\cos\theta\,A_R-\sin\theta\,A_\theta)\end{aligned}}

\begin{itemize}
\tightlist
\item
  تبدیل نمایش بردار از مختصات استوانه‌ای به کروی
\end{itemize}

\ELdisp{\vect{A}\ELbr =\uvec{r}A_r+\uvec{\phi}A_\phi+\uvec{z}A_z\quad\ELbr \ELbr \Longrightarrow\quad\ELbr \vect{A}\ELbr =\uvec{R}A_R+\uvec{\theta}A_\theta+\uvec{\phi}A_\phi}

\ELdisp{\uvec{r}\ELbr =\sin\theta\,\uvec{R}+\cos\theta\,\uvec{\theta},\qquad\ELbr  \uvec{\phi}\ELbr =\uvec{\phi},\qquad\ELbr  \uvec{z}\ELbr =\cos\theta\,\uvec{R}-\sin\theta\,\uvec{\theta}}

\ELdisp{\begin{aligned}\vect{A}&=(\sin\theta\,\uvec{R}+\cos\theta\,\uvec{\theta})A_r\\&\quad+\uvec{\phi}A_\phi\\&\quad+(\cos\theta\,\uvec{R}-\sin\theta\,\uvec{\theta})A_z\end{aligned}\quad\ELbr \ELbr \Longrightarrow\quad\ELbr \begin{aligned}\vect{A}&=\uvec{R}(\sin\theta\,A_r+\cos\theta\,A_z)\\&\quad+\uvec{\theta}(\cos\theta\,A_r-\sin\theta\,A_z)\\&\quad+\uvec{\phi}A_\phi\end{aligned}}

\end{ELunit}

\begin{ELjoin}

\begin{itemize}
\tightlist
\item
  تبدیل نمایش بردار از مختصات کروی به استوانه‌ای
\end{itemize}

\begin{ELimportant}{}

\ELdisp{\begin{bmatrix}A_r\\A_\phi\\A_z\end{bmatrix}\ELbr =\begin{bmatrix}\sin\theta&\cos\theta&0\\0&0&1\\\cos\theta&-\sin\theta&0\end{bmatrix}\begin{bmatrix}A_R\\A_\theta\\A_\phi\end{bmatrix}}

\end{ELimportant}

\end{ELjoin}

\begin{ELjoin}

\begin{itemize}
\tightlist
\item
  تبدیل نمایش بردار از مختصات استوانه‌ای به کروی
\end{itemize}

\begin{ELimportant}{}

\ELdisp{\begin{bmatrix}A_R\\A_\theta\\A_\phi\end{bmatrix}\ELbr =\begin{bmatrix}\sin\theta&0&\cos\theta\\\cos\theta&0&-\sin\theta\\0&1&0\end{bmatrix}\begin{bmatrix}A_r\\A_\phi\\A_z\end{bmatrix}}

\end{ELimportant}

\end{ELjoin}

\begin{ELunit}

\begin{itemize}
\tightlist
\item
  تبدیل نمایش بردار از مختصات کارتزین به کروی
\end{itemize}

\ELdisp{\vect{A}\ELbr =\uvec{x}A_x+\uvec{y}A_y+\uvec{z}A_z\quad\ELbr \ELbr \Longrightarrow\quad\ELbr \vect{A}\ELbr =\uvec{R}A_R+\uvec{\theta}A_\theta+\uvec{\phi}A_\phi}

\ELdisp{\left\{\begin{aligned}\uvec{x}&=\cos\phi\,\uvec{r}-\sin\phi\,\uvec{\phi}\\\uvec{y}&=\sin\phi\,\uvec{r}+\cos\phi\,\uvec{\phi}\\\uvec{z}&=\uvec{z}\end{aligned}\right.\qquad\ELbr \left\{\begin{aligned}\uvec{r}&=\sin\theta\,\uvec{R}+\cos\theta\,\uvec{\theta}\\\uvec{\phi}&=\uvec{\phi}\\\uvec{z}&=\cos\theta\,\uvec{R}-\sin\theta\,\uvec{\theta}\end{aligned}\right.}

\ELdisp{\Longrightarrow\quad\ELbr \left\{\begin{aligned}\uvec{x}&=\sin\theta\cos\phi\,\uvec{R}+\cos\theta\cos\phi\,\uvec{\theta}-\sin\phi\,\uvec{\phi}\\\uvec{y}&=\sin\theta\sin\phi\,\uvec{R}+\cos\theta\sin\phi\,\uvec{\theta}+\cos\phi\,\uvec{\phi}\\\uvec{z}&=\cos\theta\,\uvec{R}-\sin\theta\,\uvec{\theta}\end{aligned}\right.}

\end{ELunit}

\begin{ELimportant}{}

\ELdisp{\begin{bmatrix}A_R\\A_\theta\\A_\phi\end{bmatrix}\ELbr =\begin{bmatrix}\sin\theta\cos\phi&\sin\theta\sin\phi&\cos\theta\\\cos\theta\cos\phi&\cos\theta\sin\phi&-\sin\theta\\-\sin\phi&\cos\phi&0\end{bmatrix}\begin{bmatrix}A_x\\A_y\\A_z\end{bmatrix}}

\end{ELimportant}

\begin{ELunit}

\begin{itemize}
\tightlist
\item
  تبدیل نمایش بردار از مختصات کروی به کارتزین
\end{itemize}

\ELdisp{\vect{A}\ELbr =\uvec{R}A_R+\uvec{\theta}A_\theta+\uvec{\phi}A_\phi\quad\ELbr \ELbr \Longrightarrow\quad\ELbr \vect{A}\ELbr =\uvec{x}A_x+\uvec{y}A_y+\uvec{z}A_z}

\ELdisp{\left\{\begin{aligned}\uvec{R}&=\sin\theta\,\uvec{r}+\cos\theta\,\uvec{z}\\\uvec{\theta}&=\cos\theta\,\uvec{r}-\sin\theta\,\uvec{z}\\\uvec{\phi}&=\uvec{\phi}\end{aligned}\right.\qquad\ELbr \left\{\begin{aligned}\uvec{r}&=\cos\phi\,\uvec{x}+\sin\phi\,\uvec{y}\\\uvec{\phi}&=-\sin\phi\,\uvec{x}+\cos\phi\,\uvec{y}\\\uvec{z}&=\uvec{z}\end{aligned}\right.}

\ELdisp{\Longrightarrow\quad\ELbr \left\{\begin{aligned}\uvec{R}&=\sin\theta\cos\phi\,\uvec{x}+\sin\theta\sin\phi\,\uvec{y}+\cos\theta\,\uvec{z}\\\uvec{\theta}&=\cos\theta\cos\phi\,\uvec{x}+\cos\theta\sin\phi\,\uvec{y}-\sin\theta\,\uvec{z}\\\uvec{\phi}&=-\sin\phi\,\uvec{x}+\cos\phi\,\uvec{y}\end{aligned}\right.}

\end{ELunit}

\begin{ELimportant}{}

\ELdisp{\begin{bmatrix}A_x\\A_y\\A_z\end{bmatrix}\ELbr =\begin{bmatrix}\sin\theta\cos\phi&\cos\theta\cos\phi&-\sin\phi\\\sin\theta\sin\phi&\cos\theta\sin\phi&\cos\phi\\\cos\theta&-\sin\theta&0\end{bmatrix}\begin{bmatrix}A_R\\A_\theta\\A_\phi\end{bmatrix}}

\end{ELimportant}

\begin{ELunit}

\begin{itemize}
\tightlist
\item
  تبدیل از مختصات استوانه‌ای به کروی
\end{itemize}

\ELdisp{R\ELbr =\sqrt{r^2+z^2},\qquad\ELbr  \theta\ELbr =\tan^{-1}\frac{r}{z},\qquad\ELbr  \phi\ELbr =\phi}

\begin{itemize}
\tightlist
\item
  تبدیل از مختصات کروی به استوانه‌ای
\end{itemize}

\ELdisp{r\ELbr =R\sin\theta,\qquad\ELbr  \phi\ELbr =\phi,\qquad\ELbr  z\ELbr =R\cos\theta}

\begin{itemize}
\tightlist
\item
  تبدیل از مختصات کروی به کارتزین
\end{itemize}

\ELdisp{x\ELbr =R\sin\theta\cos\phi,\qquad\ELbr  y\ELbr =R\sin\theta\sin\phi,\qquad\ELbr  z\ELbr =R\cos\theta}

\begin{itemize}
\tightlist
\item
  تبدیل از مختصات کارتزین به کروی
\end{itemize}

\ELdisp{R\ELbr =\sqrt{x^2+y^2+z^2},\qquad\ELbr  \theta\ELbr =\tan^{-1}\frac{\sqrt{x^2+y^2}}{z},\qquad\ELbr  \phi\ELbr =\tan^{-1}\frac{y}{x}}

\ELfigure{m01-sph-coords}{}

\end{ELunit}

\ELhold

\ELsection{انتگرال‌های شامل توابع برداری}{انتگرال‌های شامل توابع برداری}

\begin{ELunit}

\ELfigure{m01-vector-analysis-map-calc}{}

\begin{itemize}
\tightlist
\item
  انواع انتگرال‌های شامل توابع برداری که با آن‌ها سروکار داریم
\end{itemize}

\ELdisp{\int_C\vect{F}\,d\ell\qquad\ELbr  \iint_S\vect{F}\,ds\qquad\ELbr  \iiint_V\vect{F}\,dv}
\ELdisp{\int_C V\,\dif\vect{\ell}}
\ELdisp{\int_C\vect{F}\cdot\dif\vect{\ell}\qquad\ELbr  \iint_S\vect{A}\cdot d\vect{s}}
\ELdisp{\int_C\vect{F}\times\dif\vect{\ell}}

\end{ELunit}

\begin{ELremark}{}

نکات مهم در مورد انتگرال‌گیری شامل توابع برداری

\begin{itemize}
\tightlist
\item
  استفاده از دستگاه مختصات مناسب
\item
  تعیین صحیح طول، سطح و حجم دیفرانسیلی
\item
  دقت به تغییرات بردارهای یکه \ELim{\uvec{r}}، \ELim{\uvec{R}}، \ELim{\uvec{\theta}} و \ELim{\uvec{\phi}} در بازه انتگرال‌گیری
\end{itemize}

\end{ELremark}

\begin{ELunit}

\begin{itemize}
\tightlist
\item
  انتگرال از یک تابع برداری روی منحنی، سطح و حجم مشخص
\end{itemize}

\ELdisp{\int_C\vect{F}\,d\ell\qquad\ELbr  \iint_S\vect{F}\,ds\qquad\ELbr  \iiint_V\vect{F}\,dv}

\begin{itemize}
\tightlist
\item
  حاصل یک بردار است.
\item
  کاربردهای مهم در الکترومغناطیس: محاسبه شدت میدان الکتریکی ناشی از یک توزیع بار خطی، سطحی و حجمی
\end{itemize}

\ELdisp{\left\{\begin{aligned}\vect{E}&=\frac{1}{4\pi\varepsilon_0}\int_C\frac{\rho_\ell\,d\ell'(\vect{R}-\vect{R}')}{\abs{\vect{R}-\vect{R}'}^3}\\\vect{E}&=\frac{1}{4\pi\varepsilon_0}\iint_S\frac{\rho_s\,ds'(\vect{R}-\vect{R}')}{\abs{\vect{R}-\vect{R}'}^3}\\\vect{E}&=\frac{1}{4\pi\varepsilon_0}\iiint_V\frac{\rho_v\,dv'(\vect{R}-\vect{R}')}{\abs{\vect{R}-\vect{R}'}^3}\end{aligned}\right.}

\end{ELunit}

\ELnextnum{1-9}

\begin{ELexample}{}

انتگرال زیر را روی منحنی \ELim{C} محاسبه کنید.
\ELdisp{\int_C\vect{F}\,d\ell,\qquad\ELbr  \vect{F}\ELbr =3\uvec{r}}

\ELfigure{m01-ex9}{}

\begin{ELsolution}

\ELfigure{m01-ex9-solution}{}

\ELdisp{d\ell\ELbr =r\,d\phi\ELbr =a\,d\phi}
\ELdisp{\int_C\vect{F}\,d\ell\ELbr =\int_0^\pi 3\uvec{r}\,a\,d\phi\ELbr =3a\int_0^\pi(\cos\phi\,\uvec{x}+\sin\phi\,\uvec{y})\,d\phi\ELbr =6a\uvec{y}}

\end{ELsolution}

\end{ELexample}

\begin{ELunit}

\begin{itemize}
\tightlist
\item
  انتگرال از یک تابع اسکالر روی مسیر مشخص
\end{itemize}

\ELdisp{\int_C V\,\dif\vect{\ell}}

\begin{itemize}
\tightlist
\item
  حاصل یک بردار است.
\item
  کاربردهای مهم در الکترومغناطیس: محاسبه پتانسیل مغناطیسی ناشی از یک جریان خطی
\end{itemize}

\ELdisp{\vect{A}\ELbr =\frac{\mu_0I}{4\pi}\oint_C\frac{\dif\vect{\ell}'}{\abs{\vect{R}-\vect{R}'}}}

\end{ELunit}

\ELnextnum{1-10}

\begin{ELexample}{}

انتگرال زیر را در امتداد مسیرهای زیر محاسبه کنید

\ELdisp{\int_O^P r^2\,\dif\vect{\ell},\qquad\ELbr  r^2\ELbr =x^2+y^2}

\begin{itemize}
\tightlist
\item
  مسیر \ELim{OP}
\item
  مسیر \ELim{OP_1P}
\item
  مسیر \ELim{OP_2P}
\end{itemize}

\ELfigure{m01-ex10}{}

\begin{ELsolution}

\begin{itemize}
\tightlist
\item
  در امتداد مسیر \ELim{OP}
\end{itemize}

\ELdisp{\begin{aligned}\int_O^P r^2\,\dif\vect{\ell}&=\uvec{r}\int_0^{\sqrt2}r^2\,dr=\uvec{r}\,\frac{2\sqrt2}{3}\\&=\frac{2\sqrt2}{3}(\uvec{x}\cos45^\circ+\uvec{y}\sin45^\circ)\\&=\uvec{x}\frac23+\uvec{y}\frac23\end{aligned}}

\begin{itemize}
\tightlist
\item
  در امتداد مسیر \ELim{OP_1P}
\end{itemize}

\ELdisp{\begin{aligned}\int_O^P(x^2+y^2)\,\dif\vect{\ell}&=\uvec{y}\int_O^{P_1}y^2\,dy+\uvec{x}\int_{P_1}^P(x^2+1)\,dx\\&=\uvec{y}\tfrac13y^3\Big|_0^1+\uvec{x}\left(\tfrac13x^3+x\right)\Big|_0^1\\&=\uvec{x}\frac43+\uvec{y}\frac13\end{aligned}}

\begin{itemize}
\tightlist
\item
  در امتداد مسیر \ELim{OP_2P}
\end{itemize}

\ELdisp{\begin{aligned}\int_O^P(x^2+y^2)\,\dif\vect{\ell}&=\uvec{x}\int_O^{P_2}x^2\,dx+\uvec{y}\int_{P_2}^P(1+y^2)\,dy\\&=\uvec{x}\tfrac13x^3\Big|_0^1+\uvec{y}\left(y+\tfrac13y^3\right)\Big|_0^1\\&=\uvec{x}\frac13+\uvec{y}\frac43\end{aligned}}

\end{ELsolution}

\end{ELexample}

\begin{ELunit}

\begin{itemize}
\tightlist
\item
  انتگرال خطی عددی
\end{itemize}

\ELdisp{\int_C\vect{F}\cdot\dif\vect{\ell}}

\begin{itemize}
\tightlist
\item
  حاصل یک اسکالر است.
\item
  اگر \ELim{\vect{F}} نیرو باشد

  \begin{itemize}
  \tightlist
  \item
    کار انجام شده بوسیله نیرو برای حرکت دادن جسمی از نقطه \ELim{P_1} به نقطه \ELim{P_2} در امتداد مسیر مشخص شده \ELim{C}
  \end{itemize}
\item
  کاربردهای مهم در الکترومغناطیس

  \begin{itemize}
  \item
    محاسبه اختلاف پتانسیل الکتریکی بین نقاط \ELim{P_1} و \ELim{P_2}:
    \ELdisp{V\ELbr =-\int_{P_1}^{P_2}\vect{E}\cdot\dif\vect{\ell}}
  \item
    قانون مداری آمپر:
    \ELdisp{\mu_0I\ELbr =-\oint_C\vect{B}\cdot\dif\vect{\ell}}
  \end{itemize}
\end{itemize}

\end{ELunit}

\ELnextnum{1-11}

\begin{ELexample}{}

انتگرال زیر را در مسیر ربع دایره شکل روبرو محاسبه کنید.

\ELdisp{\int_A^B\vect{F}\cdot\dif\vect{\ell},\qquad\ELbr  \vect{F}\ELbr =\uvec{x}xy-\uvec{y}2x}

\ELfigure{m01-ex11}{}

\begin{ELsolution}

\ELdisp{\begin{bmatrix}F_r\\F_\phi\\F_z\end{bmatrix}\ELbr =\begin{bmatrix}\cos\phi&\sin\phi&0\\-\sin\phi&\cos\phi&0\\0&0&1\end{bmatrix}\begin{bmatrix}xy\\-2x\\0\end{bmatrix}}
\ELdisp{\vect{F}\ELbr =\uvec{r}(xy\cos\phi-2x\sin\phi)-\uvec{\phi}(xy\sin\phi+2x\cos\phi)}
\ELdisp{\dif\vect{\ell}\ELbr =\uvec{\phi}\,3\,d\phi\qquad\ELbr  \vect{F}\cdot\dif\vect{\ell}\ELbr =-3(xy\sin\phi+2x\cos\phi)\,d\phi}
\ELdisp{x\ELbr =3\cos\phi\qquad\ELbr  y\ELbr =3\sin\phi}
\ELdisp{\int_A^B\vect{F}\cdot\dif\vect{\ell}\ELbr =\int_0^{\pi/2}-3(9\sin^2\phi\cos\phi+6\cos^2\phi)\,d\phi\ELbr =-9\left(1+\frac{\pi}{2}\right)}

\end{ELsolution}

\end{ELexample}

\begin{ELunit}

\begin{itemize}
\tightlist
\item
  انتگرال سطحی عددی
\end{itemize}

\ELdisp{\iint_S\vect{A}\cdot d\vect{s}}

\begin{itemize}
\tightlist
\item
  حاصل یک اسکالر است.
\item
  حاصل انتگرال \ELim{\ELto} شار میدان برداری \ELim{\vect{A}} که از سطح \ELim{S} می‌گذرد
\item
  سطح دیفرانسیلی برداری
\end{itemize}

\ELdisp{d\vect{s}\ELbr =\uvec{n}\,ds}

\ELfigure{m01-surface-normals}{جهت \ELim{\mathbf{a}_n} برای سطح بسته: عمود بر سطح به سمت خارج حجم؛ برای
سطح باز: وابسته به حرکت روی حاشیه سطح باز و با استفاده از قانون دست راست}

\end{ELunit}

\begin{ELjoin}

\begin{itemize}
\tightlist
\item
  کاربردهای مهم در الکترومغناطیس

  \begin{itemize}
  \item
    قانون گوس در فضای آزاد:
    \ELdisp{\oiint_S\vect{E}\cdot d\vect{s}\ELbr =\frac{Q}{\varepsilon_0}}
  \item
    محاسبه جریان الکتریکی:
    \ELdisp{\iint_S\vect{J}\cdot d\vect{s}\ELbr =I}
  \end{itemize}
\end{itemize}

\ELnextnum{1-12}

\begin{ELexample}{}

انتگرال زیر را روی سطح بسته‌ای حول محور \ELim{z} که با \ELim{z=\pm3} و \ELim{r=2} مشخص می‌شود، محاسبه کنید

\ELdisp{\oint_S\vect{F}\cdot d\vect{s},\qquad\ELbr  \vect{F}\ELbr =\uvec{r}\frac{k_1}{r}+\uvec{z}k_2z}

\ELfigure{m01-ex12}{}

\begin{ELsolution}

\ELdisp{\oint_S\vect{F}\cdot d\vect{s}\ELbr =\oint_S\vect{F}\cdot\uvec{n}\,ds\ELbr =\int_{\substack{\text{\lr{top}}\\\text{\lr{face}}}}\vect{F}\cdot\uvec{n}\,ds+\int_{\substack{\text{\lr{bottom}}\\\text{\lr{face}}}}\vect{F}\cdot\uvec{n}\,ds+\int_{\substack{\text{\lr{side}}\\\text{\lr{wall}}}}\vect{F}\cdot\uvec{n}\,ds}

\begin{itemize}
\tightlist
\item
  سطح بالایی
\end{itemize}

\ELdisp{z\ELbr =3,\quad\ELbr  \uvec{n}\ELbr =\uvec{z},\quad\ELbr  \vect{F}\cdot\uvec{n}\ELbr =k_2z\ELbr =3k_2,\quad\ELbr  ds\ELbr =r\,dr\,d\phi}
\ELdisp{\int_{\substack{\text{\lr{top}}\\\text{\lr{face}}}}\vect{F}\cdot\uvec{n}\,ds\ELbr =\int_0^{2\pi}\!\!\int_0^2 3k_2\,r\,dr\,d\phi\ELbr =12\pi k_2}

\begin{itemize}
\tightlist
\item
  سطح پایینی
\end{itemize}

\ELdisp{z\ELbr =-3,\quad\ELbr  \uvec{n}\ELbr =-\uvec{z},\quad\ELbr  \vect{F}\cdot\uvec{n}\ELbr =-k_2z\ELbr =3k_2,\quad\ELbr  ds\ELbr =r\,dr\,d\phi}
\ELdisp{\int_{\substack{\text{\lr{bottom}}\\\text{\lr{face}}}}\vect{F}\cdot\uvec{n}\,ds\ELbr =12\pi k_2}

\begin{itemize}
\tightlist
\item
  دیواره جانبی
\end{itemize}

\ELdisp{r\ELbr =2,\quad\ELbr  \uvec{n}\ELbr =\uvec{r},\quad\ELbr  \vect{F}\cdot\uvec{n}\ELbr =\frac{k_1}{r}\ELbr =\frac{k_1}{2},\quad\ELbr  ds\ELbr =r\,d\phi\,dz\ELbr =2\,d\phi\,dz}
\ELdisp{\int_{\substack{\text{\lr{side}}\\\text{\lr{wall}}}}\vect{F}\cdot\uvec{n}\,ds\ELbr =\int_{-3}^3\!\int_0^{2\pi}k_1\,d\phi\,dz\ELbr =12\pi k_1}

\ELdisp{\oint_S\vect{F}\cdot d\vect{s}\ELbr =12\pi k_2+12\pi k_2+12\pi k_1\ELbr =12\pi(k_1+2k_2)}

\end{ELsolution}

\end{ELexample}

\end{ELjoin}

\begin{ELunit}

\begin{itemize}
\tightlist
\item
  انتگرال خطی برداری
\end{itemize}

\ELdisp{\int_C\vect{F}\times\dif\vect{\ell}}

\end{ELunit}

\begin{ELjoin}

\begin{itemize}
\tightlist
\item
  حاصل یک بردار است.
\item
  کاربردهای مهم در الکترومغناطیس

  \begin{itemize}
  \item
    قانون بیوساوار:
    \ELdisp{\vect{B}\ELbr =\frac{\mu_0I}{4\pi}\oint_C\frac{\dif\vect{\ell}\times(\vect{R}-\vect{R}')}{\abs{\vect{R}-\vect{R}'}^3}}
  \item
    محاسبه نیروی مغناطیسی وارد بر سیم حامل جریان در میدان مغناطیسی:
    \ELdisp{\vect{F}\ELbr =I\oint_C\dif\vect{\ell}\times\vect{B}}
  \end{itemize}
\end{itemize}

\ELnextnum{1-13}

\begin{ELexample}{}

انتگرال زیر را در مسیر \ELim{C} محاسبه کنید.

\ELdisp{\int_C\vect{F}\times\dif\vect{\ell},\qquad\ELbr  \vect{F}\ELbr =2\uvec{z}}

\ELfigure{m01-ex13}{}

\begin{ELsolution}

\ELdisp{\dif\vect{\ell}\ELbr =r\,d\phi\,\uvec{\phi}\ELbr =a\,d\phi\,\uvec{\phi}}
\ELdisp{\begin{aligned}\int_C\vect{F}\times\dif\vect{\ell}&=\int_{2\pi}^0 2\uvec{z}\times a\,d\phi\,\uvec{\phi}=-2a\int_{2\pi}^0\uvec{r}\,d\phi\\&=-2a\int_{2\pi}^0(\cos\phi\,\uvec{x}+\sin\phi\,\uvec{y})\,d\phi=0\end{aligned}}

\end{ELsolution}

\end{ELexample}

\end{ELjoin}

\ELhold

\ELsection{عملیات مشتق‌گیری - گرادیان}{عملیات مشتق‌گیری - گرادیان}

\begin{ELunit}

\begin{itemize}
\tightlist
\item
  \textbf{مفاهیم پیش‌نیاز برای ورود به بحث گرادیان}

  \begin{itemize}
  \tightlist
  \item
    مفهوم مشتق از یک تابع تک‌متغیره

    \begin{itemize}
    \tightlist
    \item
      نرخ تغییرات تابع در نقطه مورد نظر
    \end{itemize}
  \item
    مفهوم مشتق جهت‌دار از یک تابع چندمتغیره

    \begin{itemize}
    \tightlist
    \item
      نرخ تغییرات تابع در نقطه مورد نظر و در یک جهت مشخص شده
    \end{itemize}
  \end{itemize}
\end{itemize}

\ELfigure{m01-derivative-concepts}{}

\end{ELunit}

\begin{ELdefinition}{گرادیان}

گرادیانِ یک \textbf{کمیت اسکالر}، یک \textbf{کمیت برداری} است که در جهتی قرار می‌گیرد که کمیت اسکالر بیشترین افزایش را دارد و اندازه آن برابر با حداکثر نرخ افزایش کمیت اسکالر است.

\ELdisp{\operatorname{grad}v\ELbr =\nabla v\ELbr =\uvec{n}\frac{dv}{dn}}

\end{ELdefinition}

\begin{ELunit}

\begin{itemize}
\tightlist
\item
  مشتق جهت‌دار \ELim{v} در جهت \ELim{\dif\vect{\ell}}:
\end{itemize}

\ELdisp{\frac{dv}{d\ell}\ELbr =\nabla v\cdot\uvec{\ell}}

\begin{itemize}
\tightlist
\item
  گرادیان \ELim{v} بر سطوح \ELim{v} ثابت عمود است.
\end{itemize}

\ELfigure{m01-gradient-map}{}

\end{ELunit}

\begin{ELjoin}

\begin{itemize}
\tightlist
\item
  رابطه گرادیان در مختصات کارتزین به صورت زیر است
\end{itemize}

\begin{ELimportant}{}

\ELdisp{\nabla v\ELbr =\frac{\partial v}{\partial x}\uvec{x}+\frac{\partial v}{\partial y}\uvec{y}+\frac{\partial v}{\partial z}\uvec{z}}

\end{ELimportant}

\end{ELjoin}

\begin{ELjoin}

\begin{itemize}
\tightlist
\item
  رابطه کلی گرادیان
\end{itemize}

\begin{ELimportant}{}

\ELdisp{\nabla V\ELbr =\uvec{u_1}\frac{\partial V}{h_1\,\partial u_1}+\uvec{u_2}\frac{\partial V}{h_2\,\partial u_2}+\uvec{u_3}\frac{\partial V}{h_3\,\partial u_3}}

\end{ELimportant}

\end{ELjoin}

\Needspace{10\baselineskip}
\begin{longtable}[c]{llll}
\toprule\noalign{}
\ELtablehead ضرایب متری & کارتزین & استوانه‌ای & کروی \\\noalign{\vskip 4pt}
\midrule\noalign{}
\endhead
\bottomrule\noalign{}
\endlastfoot
\ELim{h_1} & \ELim{1} & \ELim{1} & \ELim{1} \\\noalign{\vskip 4pt}
\ELim{h_2} & \ELim{1} & \ELim{r} & \ELim{R} \\\noalign{\vskip 4pt}
\ELim{h_3} & \ELim{1} & \ELim{1} & \ELim{R\sin\theta} \\\noalign{\vskip 4pt}
\end{longtable}

\begin{ELunit}

\begin{itemize}
\tightlist
\item
  \textbf{مفهوم گرادیان}

  \begin{itemize}
  \tightlist
  \item
    نمایش یک تپه با کانتورهای ارتفاع \ELim{\ELto} \ELim{h(x,y)}
  \item
    اگر توپی در نقطه \ELim{P} قرار داشته باشد، در چه جهتی حرکت می‌کند و مقدار نیروی وارد بر آن چقدر است؟

    \begin{itemize}
    \tightlist
    \item
      پاسخ متناسب با \ELim{-\nabla h} است.
    \end{itemize}
  \end{itemize}
\end{itemize}

\ELfigure{m01-gradient-hill}{}

\end{ELunit}

\ELnextnum{1-14}

\begin{ELexample}{}

شدت میدان الکتریکی \ELim{\vect{E}} به صورت منهای گرادیان پتانسیل الکتریکی عددی \ELim{V} قابل دستیابی است. \ELim{\vect{E}} را در نقطه \ELim{(1,1,0)} بیابید اگر
\ELdisp{V\ELbr =E_0R\cos\theta}

\begin{ELsolution}

\ELdisp{\vect{E}\ELbr =-\nabla V}
\ELdisp{\begin{aligned}\vect{E}&=-\left[\uvec{R}\frac{\partial}{\partial R}+\uvec{\theta}\frac{\partial}{R\,\partial\theta}+\uvec{\phi}\frac{\partial}{R\sin\theta\,\partial\phi}\right]E_0R\cos\theta\\&=-(\uvec{R}\cos\theta-\uvec{\theta}\sin\theta)E_0\end{aligned}}
\ELdisp{\vect{E}\ELbr =-\uvec{z}E_0}

\end{ELsolution}

\end{ELexample}

\ELhold

\ELsection{عملیات مشتق‌گیری - دیورژانس}{عملیات مشتق‌گیری - دیورژانس}

\begin{ELunit}

\begin{itemize}
\tightlist
\item
  \textbf{مفاهیم پیش‌نیاز برای ورود به بحث دیورژانس}

  \begin{itemize}
  \tightlist
  \item
    در مطالعه میدان‌های برداری، راحت‌تر است که تغییرات میدان را به صورت ترسیمی توسط خطوط شار نمایش دهیم. این‌ها خطوط یا منحنی‌های جهت‌داری هستند که در هر نقطه جهت میدان برداری را مشخص می‌کنند.
  \end{itemize}
\end{itemize}

\ELfigure{m01-flux-lines}{}

\begin{itemize}
\tightlist
\item
  شار یک میدان برداری مشابه جریان یک سیال مانند آب است.
\item
  اگر \ELim{\vect{A}} یک بردار چگالی شار باشد، کل شار خروجی از سطح \ELim{S} برابر است با
\end{itemize}

\ELdisp{\oint_S\vect{A}\cdot d\vect{s}}

\begin{itemize}
\tightlist
\item
  در حجمی با یک سطح بسته، تنها وقتی شار اضافی ورودی یا خروجی وجود خواهد داشت که این حجم به ترتیب دارای یک چاه یا یک منبع باشد.

  \begin{itemize}
  \tightlist
  \item
    شار خالص خروجی از \ELim{v_a} مثبت است: وجود منبع در حجم \ELim{v_a}
  \item
    شار خالص خروجی از \ELim{v_b} منفی است: وجود چاه در حجم \ELim{v_b}
  \item
    شار خالص خروجی از \ELim{v_c} صفر است.
  \end{itemize}
\end{itemize}

\ELfigure{m01-source-sink}{}

\end{ELunit}

\begin{ELdefinition}{دیورژانس}

دیورژانس \textbf{میدان برداری} \ELim{\vect{A}} در یک نقطه عبارتست از شار خالص خروجی \ELim{\vect{A}} در واحد حجم وقتی که این حجم حول نقطه به سمت صفر میل می‌کند

\ELdisp{\operatorname{div}\vect{A}\triangleq\lim_{\Delta v\to0}\frac{\oint_S\vect{A}\cdot d\vect{s}}{\Delta v}}

\end{ELdefinition}

\begin{ELjoin}

\begin{itemize}
\tightlist
\item
  در مختصات کارتزین رابطه دیورژانس به صورت زیر است
\end{itemize}

\begin{ELimportant}{}

\ELdisp{\operatorname{div}\vect{A}\ELbr =\frac{\partial A_x}{\partial x}+\frac{\partial A_y}{\partial y}+\frac{\partial A_z}{\partial z}\qquad\ELbr \qquad\ELbr  \nabla\cdot\vect{A}\equiv\operatorname{div}\vect{A}}

\end{ELimportant}

\end{ELjoin}

\begin{ELjoin}

\begin{itemize}
\tightlist
\item
  رابطه کلی دیورژانس
\end{itemize}

\begin{ELimportant}{}

\ELdisp{\nabla\cdot\vect{A}\ELbr =\frac{1}{h_1h_2h_3}\left[\frac{\partial}{\partial u_1}(h_2h_3A_1)+\frac{\partial}{\partial u_2}(h_1h_3A_2)+\frac{\partial}{\partial u_3}(h_1h_2A_3)\right]}

\end{ELimportant}

\end{ELjoin}

\ELnextnum{1-15}

\begin{ELexample}{}

چگالی شار مغناطیسی \ELim{\vect{B}} در بیرون یک سیم طویل حامل جریان به صورت دایره‌ای و متناسب با معکوس فاصله از محور سیم است. دیورژانس \ELim{\vect{B}} را بیابید.

\begin{ELsolution}

\begin{itemize}
\tightlist
\item
  با فرض منطبق بودن سیم بر محور \ELim{z}
\end{itemize}

\ELdisp{\vect{B}\ELbr =\uvec{\phi}\frac{k}{r}}
\ELdisp{\nabla\cdot\vect{B}\ELbr =\frac1r\frac{\partial}{\partial r}(rB_r)+\frac1r\frac{\partial B_\phi}{\partial\phi}+\frac{\partial B_z}{\partial z}}
\ELdisp{\nabla\cdot\vect{B}\ELbr =0}

\ELfigure{m01-ex15}{}

\end{ELsolution}

\end{ELexample}

\ELhold

\ELsection{عملیات مشتق‌گیری - کرل}{عملیات مشتق‌گیری - کرل}

\begin{ELunit}

\begin{itemize}
\tightlist
\item
  \textbf{مفاهیم پیش‌نیاز برای ورود به بحث کرل}

  \begin{itemize}
  \tightlist
  \item
    گردش یک میدان برداری به دور یک مسیر بسته
  \end{itemize}
\end{itemize}

\ELdisp{\text{\lr{Circulation of }}\vect{A}\text{\lr{ around contour }}C\triangleq\oint_C\vect{A}\cdot\dif\vect{\ell}}

\begin{itemize}
\tightlist
\item
  معنای فیزیکی گردش به نوع میدانی که بردار \ELim{\vect{A}} نمایش می‌دهد بستگی دارد.

  \begin{itemize}
  \tightlist
  \item
    اگر \ELim{\vect{A}} نیروی وارد بر یک جسم باشد، گردش \ELim{\vect{A}} کار انجام شده توسط نیرو برای حرکت جسم به دور مسیر است.
  \end{itemize}
\end{itemize}

\end{ELunit}

\begin{ELdefinition}{کرل}

کرل \textbf{میدان برداری} \ELim{\vect{A}} در یک نقطه \textbf{برداری} است که اندازه آن حداکثر گردش خالص \ELim{\vect{A}} در واحد سطح است وقتی که سطح حول نقطه به سمت صفر میل می‌کند و جهت آن جهت عمود سطح است زمانی که سطح طوری جهت داده شده باشد که گردش خالص را حداکثر کند.

\ELdisp{\operatorname{curl}\vect{A}\equiv\nabla\times\vect{A}\triangleq\lim_{\Delta s\to0}\frac{1}{\Delta s}\left[\uvec{n}\oint_C\vect{A}\cdot\dif\vect{\ell}\right]_{\max}}

\end{ELdefinition}

\begin{ELjoin}

\begin{itemize}
\tightlist
\item
  در مختصات کارتزین رابطه کرل به صورت زیر است
\end{itemize}

\begin{ELimportant}{}

\ELdisp{\nabla\times\vect{A}\ELbr =\begin{vmatrix}\uvec{x}&\uvec{y}&\uvec{z}\\[1mm]\dfrac{\partial}{\partial x}&\dfrac{\partial}{\partial y}&\dfrac{\partial}{\partial z}\\[3mm]A_x&A_y&A_z\end{vmatrix}}

\end{ELimportant}

\end{ELjoin}

\begin{ELjoin}

\begin{itemize}
\tightlist
\item
  رابطه کلی کرل
\end{itemize}

\begin{ELimportant}{}

\ELdisp{\nabla\times\vect{A}\ELbr =\frac{1}{h_1h_2h_3}\begin{vmatrix}\uvec{u_1}h_1&\uvec{u_2}h_2&\uvec{u_3}h_3\\[1mm]\dfrac{\partial}{\partial u_1}&\dfrac{\partial}{\partial u_2}&\dfrac{\partial}{\partial u_3}\\[3mm]h_1A_1&h_2A_2&h_3A_3\end{vmatrix}}

\end{ELimportant}

\end{ELjoin}

\begin{ELunit}

\begin{itemize}
\tightlist
\item
  \textbf{مفهوم کرل}

  \begin{itemize}
  \tightlist
  \item
    چرخش یک چرخ پره‌دار درون جریان آب با بردار سرعت \ELim{u}

    \begin{itemize}
    \tightlist
    \item
      حداکثر سرعت زاویه‌ای چرخ پره‌دار متناسب با اندازه کرل بردار سرعت است.
    \item
      محور چرخ با استفاده از قانون دست راست در جهت کرل قرار می‌گیرد.
    \end{itemize}
  \end{itemize}
\end{itemize}

\ELfigure{m01-paddle-wheel}{}

\end{ELunit}

\ELnextnum{1-16}

\begin{ELexample}{}

کرل بردار زیر را محاسبه کنید
\ELdisp{\vect{A}\ELbr =\uvec{\phi}\left(\frac{k}{r}\right)}

\begin{ELsolution}

\ELdisp{\nabla\times\vect{A}\ELbr =\frac1r\begin{vmatrix}\uvec{r}&\uvec{\phi}r&\uvec{z}\\[1mm]\dfrac{\partial}{\partial r}&\dfrac{\partial}{\partial\phi}&\dfrac{\partial}{\partial z}\\[3mm]0&k&0\end{vmatrix}\ELbr =0}

\ELfigure{m01-ex16}{}

\end{ELsolution}

\end{ELexample}

\ELhold

\ELsection{قضایای مهم برداری}{قضایای مهم برداری}

\begin{ELunit}

\ELfigure{m01-vector-analysis-map-thm}{}

\end{ELunit}

\begin{ELtheorem}{قضیه دیورژانس}

انتگرال حجمی دیورژانس یک میدان برداری با شار کل خروجی بردار از سطح در برگیرنده حجم برابر است.
\ELdisp{\int_V\nabla\cdot\vect{A}\,dv\ELbr =\oint_S\vect{A}\cdot d\vect{s}}

\end{ELtheorem}

\begin{ELunit}

\ELfigure{m01-divergence-theorem}{}

\end{ELunit}

\begin{ELtheorem}{قضیه استوکس}

انتگرال سطحی کرل یک میدان برداری، روی یک سطح باز برابر انتگرال خطی بسته بردار روی مسیری است که سطح را در بر می‌گیرد.
\ELdisp{\int_S(\nabla\times\vect{A})\cdot d\vect{s}\ELbr =\oint_C\vect{A}\cdot\dif\vect{\ell}}

\end{ELtheorem}

\begin{ELunit}

\ELfigure{m01-stokes-theorem}{}

\end{ELunit}

\begin{ELtheorem}{قضیه هلم‌هولتز}

یک میدان برداری تا حد یک ثابت افزودنی تعیین می‌شود اگر هم دیورژانس و هم کرل آن مشخص باشد.

\end{ELtheorem}

\ELnextnum{1-17}

\begin{ELexample}{}

اعتبار قضیه دیورژانس را در مورد ناحیه پوسته‌ای محصور به سطوح کروی \ELim{R=R_1} و \ELim{R=R_2} \ELim{(R_2>R_1)} به مرکز مبدأ مختصات برای میدان برداری زیر تحقیق نمایید.
\ELdisp{\vect{F}\ELbr =\uvec{R}kR}

\ELfigure{m01-ex17}{}

\begin{ELsolution}

\ELdisp{\int_V\nabla\cdot\vect{A}\,dv\ELbr =\oint_S\vect{A}\cdot d\vect{s}}

\begin{itemize}
\tightlist
\item
  برای سطح بیرونی
\end{itemize}

\ELdisp{R\ELbr =R_2,\qquad\ELbr  d\vect{s}\ELbr =\uvec{R}R_2^2\sin\theta\,d\theta\,d\phi}
\ELdisp{\int_{\substack{\text{\lr{outer}}\\\text{\lr{surface}}}}\vect{F}\cdot d\vect{s}\ELbr =\int_0^{2\pi}\!\!\int_0^\pi(kR_2)R_2^2\sin\theta\,d\theta\,d\phi\ELbr =4\pi kR_2^3}

\begin{itemize}
\tightlist
\item
  برای سطح داخلی
\end{itemize}

\ELdisp{R\ELbr =R_1,\qquad\ELbr  d\vect{s}\ELbr =-\uvec{R}R_1^2\sin\theta\,d\theta\,d\phi}
\ELdisp{\int_{\substack{\text{\lr{inner}}\\\text{\lr{surface}}}}\vect{F}\cdot d\vect{s}\ELbr =-\int_0^{2\pi}\!\!\int_0^\pi(kR_1)R_1^2\sin\theta\,d\theta\,d\phi\ELbr =-4\pi kR_1^3}

\ELdisp{\oint_S\vect{F}\cdot d\vect{s}\ELbr =4\pi k(R_2^3-R_1^3)}
\ELdisp{\nabla\cdot\vect{F}\ELbr =\frac{1}{R^2}\frac{\partial}{\partial R}(R^2F_R)\ELbr =\frac{1}{R^2}\frac{\partial}{\partial R}(kR^3)\ELbr =3k}
\ELdisp{\int_V\nabla\cdot\vect{F}\,dv\ELbr =(\nabla\cdot\vect{F})V\ELbr =4\pi k(R_2^3-R_1^3)}

\end{ELsolution}

\end{ELexample}

\ELnextnum{1-18}

\begin{ELexample}{}

اعتبار قضیه استوکس را روی یک چهارم قرص مدوری به شعاع ۳ در ربع اول برای میدان برداری زیر تحقیق نمایید.
\ELdisp{\vect{F}\ELbr =\uvec{x}xy-\uvec{y}2x}

\ELfigure{m01-ex18}{}

\begin{ELsolution}

\ELdisp{\int_S(\nabla\times\vect{A})\cdot d\vect{s}\ELbr =\oint_C\vect{A}\cdot\dif\vect{\ell}}
\ELdisp{\nabla\times\vect{F}\ELbr =\begin{vmatrix}\uvec{x}&\uvec{y}&\uvec{z}\\[1mm]\dfrac{\partial}{\partial x}&\dfrac{\partial}{\partial y}&\dfrac{\partial}{\partial z}\\[3mm]xy&-2x&0\end{vmatrix}\ELbr =-\uvec{z}(2+x)}
\ELdisp{\int_S(\nabla\times\vect{F})\cdot d\vect{s}\ELbr =\int_0^{\pi/2}\!\!\int_0^3-\uvec{z}(2+r\cos\phi)\cdot r\,dr\,d\phi\,\uvec{z}\ELbr =-9\left(1+\frac{\pi}{2}\right)}

\begin{itemize}
\tightlist
\item
  از \ELim{B} تا \ELim{O}
\end{itemize}

\ELdisp{x\ELbr =0,\ \text{\lr{and}}\ \vect{F}\cdot\dif\vect{\ell}\ELbr =\vect{F}\cdot(\uvec{y}\,dy)\ELbr =2x\,dy\ELbr =0}

\begin{itemize}
\tightlist
\item
  از \ELim{O} تا \ELim{A}
\end{itemize}

\ELdisp{y\ELbr =0,\ \text{\lr{and}}\ \vect{F}\cdot\dif\vect{\ell}\ELbr =\vect{F}\cdot(\uvec{x}\,dx)\ELbr =xy\,dx\ELbr =0}

\ELdisp{\oint_{ABOA}\vect{F}\cdot\dif\vect{\ell}\ELbr =\int_A^B\vect{F}\cdot\dif\vect{\ell}\ELbr =-9\left(1+\frac{\pi}{2}\right)}

(در مثال \ELnumc{1-11} بررسی شده بود)

\end{ELsolution}

\end{ELexample}

\ELhold

\ELsection{اتحادهای مهم برداری}{اتحادهای مهم برداری}

\begin{ELimportant}{اتحاد \lr{I}}

کرل گرادیان هر کمیت عددی برابر با صفر است
\ELdisp{\nabla\times(\nabla V)\equiv0}

\end{ELimportant}

\begin{ELunit}

\begin{itemize}
\tightlist
\item
  یک میدان برداری بدون کرل را یک میدان \textbf{غیرگردشی} یا \textbf{ابقایی} گویند.
\item
  اگر یک کمیت برداری بدون کرل باشد، آنگاه می‌توان آن را به صورت گرادیان یک کمیت عددی بیان کرد.
\end{itemize}

\end{ELunit}

\begin{ELimportant}{اتحاد \lr{II}}

دیورژانس کرل هر میدان برداری برابر با صفر است.
\ELdisp{\nabla\cdot(\nabla\times\vect{A})\equiv0}

\end{ELimportant}

\begin{ELunit}

\begin{itemize}
\tightlist
\item
  یک میدان برداری بدون دیورژانس را یک میدان \textbf{سلونوئیدی} گویند.
\item
  اگر یک کمیت برداری بدون دیورژانس باشد، آنگاه می‌توان آن را به صورت کرل یک کمیت برداری بیان کرد.
\end{itemize}

\end{ELunit}
